
\subsubsection{Execution-Based Code Generation Benchmarks}

\textbf{HumanEval} (Chen et al., 2021) provides 164 hand-crafted Python programming problems designed to test algorithmic problem-solving. Each problem includes a function signature, docstring specification, and unit tests. The benchmark reports pass@k metrics (percentage of problems solved with k samples).

\textbf{MBPP} (Austin et al., 2021) contains 974 Python programming problems from introductory programming exercises. The benchmark emphasizes practical programming patterns representing common coding scenarios. Problems include string manipulation, data structure operations, and basic algorithms.

\textbf{APPS} (Hendrycks et al., 2021) includes 10,000 competitive programming problems from coding competition platforms. Problems are stratified by difficulty (introductory, interview, competition levels) and require complex algorithmic reasoning.

These benchmark papers document different design philosophies but do not test whether design differences produce independent model rankings.

\subsubsection{Multi-Benchmark Evaluation}

The BigCode Evaluation Harness provides infrastructure for standardized evaluation across code generation benchmarks. It reports pass@k scores independently for each benchmark without analyzing structural relationships between benchmarks. Code LLM leaderboards aggregate multi-benchmark performance but treat each benchmark as an independent evaluation axis.

Current multi-benchmark frameworks assume dimensional independence without empirical validation. They report HumanEval, MBPP, and APPS scores as separate columns, implicitly treating them as measuring distinct competencies.

\subsubsection{Factor Analysis in Evaluation}

Factor analysis is used in psychometrics to validate whether multiple measurements capture independent dimensions or measure a shared construct. In educational testing, confirmatory factor analysis determines whether reading comprehension, mathematical reasoning, and spatial reasoning represent independent cognitive abilities or load on a general intelligence factor.

Factor-analytic methods have limited application to machine learning benchmark validation. Benchmark design typically emphasizes face validity rather than empirical validation of dimensional structure.

Our work differs by empirically testing ranking correlation between benchmarks rather than assuming independence from design philosophy differences.
